%-------------------------
% Resume in Latex
% Date December 2025
% Author : John O'Shea
%------------------------

\documentclass[letterpaper,10pt]{article}

\usepackage{makecell}
\usepackage[link=off]{phonenumbers}

\usepackage{latexsym}
\usepackage[empty]{fullpage}
\usepackage{titlesec}
\usepackage{marvosym}
\usepackage[usenames,dvipsnames]{color}
\usepackage{verbatim}
\usepackage{enumitem}
\usepackage[pdftex]{hyperref}
\usepackage{fancyhdr}
\usepackage{graphicx}



\pagestyle{fancy}
\fancyhf{} % clear all header and footer fields
\fancyfoot{}
\renewcommand{\headrulewidth}{0pt}
\renewcommand{\footrulewidth}{0pt}
\usepackage[margin=0.3in]{geometry}
% Adjust margins
\addtolength{\oddsidemargin}{-0.0in}
\addtolength{\evensidemargin}{-0.0in}
\addtolength{\textwidth}{0in}
\addtolength{\topmargin}{20pt}
\addtolength{\textheight}{0.0in}

\urlstyle{same}

\usepackage{xcolor}% http://ctan.org/pkg/xcolor
\usepackage{hyperref}% http://ctan.org/pkg/hyperref
\hypersetup{
  colorlinks=true,
  linkcolor=blue!50!red,
  linkbordercolor=red,
  urlcolor=blue!70!black
}

\raggedbottom
\raggedright
\setlength{\tabcolsep}{0in}

% Sections formatting
\titleformat{\section}{
  \vspace{-10pt}\scshape\raggedright\large
}{}{0em}{}[\color{black}\titlerule \vspace{-7pt}]

%-------------------------
% Custom commands
\def \ifempty#1{\def\temp{#1} \ifx\temp\empty }

\newcommand{\resumeItem}[2]{
  \item\small{
  	\ifempty{#1}#2\else\textbf{#1}{: #2 \vspace{-2pt}}\fi
  }
}

\usepackage[dvipsnames]{xcolor}
\definecolor{mygray}{gray}{0}
\usepackage{fancybox}

\usepackage{lmodern}
\usepackage{tikz}

% Style definition
\tikzset{rndblock/.style={rounded corners,rectangle,draw,outer sep=0pt}}

% Command Definition
% 1 optional to customize the aspect, 2 mandatory: text to be framed
\newcommand{\tframed}[2][]{\tikz[baseline=(h.base)]\node[rndblock,#1] (h) {\color{black}{#2}};}

\newcommand*{\mystrut}{\rule[-0.2\baselineskip]{0pt}{0.8\baselineskip}}
\newcommand{\skill}[1]{\tframed[lightgray]{\mystrut#1}}


\newcommand{\resumeSubheading}[4]{
  \vspace{-1pt}\item
    \begin{tabular*}{0.97\textwidth}{l@{\extracolsep{\fill}}r}
      \textbf{#1} & \textcolor{mygray}{#2} \\
      \textit{\small#3} & \textcolor{mygray}{\textit{\small #4}} \\
    \end{tabular*}\vspace{-5pt}
}

\newcommand{\resumeExpSubheading}[5]{
  \vspace{-1pt}\item
    \begin{tabular*}{0.97\textwidth}{l@{\extracolsep{\fill}}r}
      \textbf{#1}  & \textcolor{mygray}{#2} \\
      \textit{\small#3} & \textcolor{mygray}{\textit{\small #4}} \\
      {\scriptsize#5}
    \end{tabular*}\vspace{4pt}
}

\newcommand{\resumeProjSubheading}[4]{
  \vspace{-1pt}\item
    \begin{tabular*}{0.97\textwidth}{l@{\extracolsep{\fill}}r}
      \textbf{#1}  & \textcolor{mygray}{#2} \\
      \scriptsize {#3} & \textcolor{mygray}{\textit{\small #4}} \\
    \end{tabular*}\vspace{4pt}
}

\newcommand{\resumeSubItem}[2]{\resumeItem{#1}{#2}\vspace{-4pt}}

\renewcommand{\labelitemii}{$\circ$}

\newcommand{\resumeSubHeadingListStart}{\begin{itemize}[leftmargin=*]}
\newcommand{\resumeSubHeadingListEnd}{\end{itemize}}
\newcommand{\resumeItemListStart}{\begin{itemize}[leftmargin=0.2in]}
\newcommand{\resumeItemListEnd}{\end{itemize}\vspace{-5pt}}

\usepackage{changepage}
\newcommand{\resumeDesc}[1]{\begin{adjustwidth}{5pt}{0pt}\vspace{-2pt}{\small{#1}}\end{adjustwidth}}

%-------------------------------------------
%%%%%%  CV STARTS HERE  %%%%%%%%%%%%%%%%%%%%%%%%%%%%


\begin{document}

\begin{tabular*}{\textwidth}{l@{\extracolsep{\fill}}r}
  \includegraphics[width=0.1\linewidth]{photo-joshea.png} \\
  \textbf{\Large John O'Shea}   
  \\Email : \href{mailto:oshea.john@gmail.com}{oshea.john@gmail.com}, Github: \href{https://github.com/jfoshea}{jfoshea}, LinkedIn : \href{https://www.linkedin.com/in/john-oshea/}{john-oshea} \\
  
\end{tabular*}


\section{Summary}
    \resumeDesc
    
    {Experienced engineer and self‑directed learner with a strong record of innovation across the high‑tech industry. I bring hands‑on expertise across both software and hardware development, with extensive work in server and data‑storage technologies including sRIO, InfiniBand, and PCIe on Linux and bare‑metal embedded platforms. I have also knowledge and delivered in areas such Root of Trust, OS level projects, GPU software and sensors/control in ADAS. My background spans advanced platform engineering and complex, multi‑layered technology stacks, where I excel at systems‑level thinking, cross‑domain problem‑solving, and architecting robust, scalable, and high‑performance platforms. I am the inventor on five granted U.S. patents (7987229, 8156220, 8090789, 8645623, 7631128) and hold U.S. citizenship.}
    
%-----------EXPERIENCE-----------------
\section{Experience}
  \resumeSubHeadingListStart
      \resumeExpSubheading
      {Qualcomm }{Cork, Ireland}
      {\makecell[l]{Director of Engineering / Principal Engineer, GPU Software, Core Platform Software for Snapdragon SoC \\ Senior Staff Engineer/Manager, Root of Trust IP development/validation for Snapdragon SoC}}{\makecell[r]{Nov. 2022 - Present \\ Nov. 2019 - Aug. 2020}}
      {\skill{Leadership} \skill{C, C++} \skill{GPU SW, HW Architecture} \skill{Root Of Trust Security}}
      \resumeDesc{\begin{itemize}
          \item Built and led high‑performance software engineering teams delivering GPU software for Snapdragon Adreno, spanning Vulkan, OpenCL, SYCL, DirectX 12, research initiatives, and internal developer tools.
        \item Directed software engineering teams for AI data‑center inference products, with a focus on core platform software layer, including SoC firmware, peripheral interfaces, bootloaders, hypervisors, and power‑management systems.
        \item Led the software validation team for a new Hardware Root of Trust (RoT) IP within the Snapdragon SoC. Designed and implemented C, Python, and CMM‑based validation frameworks to exercise and verify cryptographic hardware blocks and multiple RISC‑V and custom CPU cores inside the RoT subsystem..
      \end{itemize}}
      \resumeExpSubheading
      {Helgen Technologies}{Cork, Ireland and Baltimore, MD}
      {Cofounder, Engineering Lead and Technical consultant}{Sept. 2020 - Dec. 2022}
      {\skill{Leadership} \skill{Proof of concepts} \skill{Prototyping} \skill{Customer Focus}}
      \resumeDesc{\begin{itemize}
          \item Co-Founder of a startup focused on autonomous systems, where I built and led a small engineering team developing prototype automation solutions.
          \item Boot strapped the company by working as sub-contractor for Automotive sensor startup and Drone POC initiatives. Helped develop C and C++ code for embedded development.
      \end{itemize}}

      \resumeExpSubheading
      {Jaguar Land Rover}{Shannon, Ireland}
      {Senior Software Engineer, Software Architect}{Nov. 2018 - Nov. 2019 }      
      {\skill{C++} \skill{ADAS} \skill{SW Architecture}}
      \resumeDesc{\begin{itemize}
          \item Developed vehicle motion controller algorithms for use on ADAS Level 3/4 capable self-driving software stack. Developed C++14 code for ADAS prototyping and development.
          \item Software Architecture for next-gen electrical vehicle architecture platform.
      \end{itemize}}

      \resumeExpSubheading
      {Dell Technologies, EMC Corporation}{Hopkinton, MA}
      {\makecell[l]{Principal Software Engineer, Platform SW team \\ Principal Firmware/Hardware Engineer, Platform HW team \\ Senior Hardware Engineer, Platform HW team}}
      {\makecell[r]{Jan. 2012 - Oct. 2018 \\ Jan. 2006 - Jan. 2012 \\ Jan. 2000 - Jan. 2006}}
      {\skill{C, C++, Python} \skill{Linux, U-Boot, Device Tree, Buildroot, UEFI} \skill{BMC,IPMI} \skill{sRIO, Infiniband, Ethernet,PCIe}}
      \resumeDesc{\begin{itemize}
          \item Developed and maintained Embedded Linux software for data‑center hardware platforms, including U‑Boot, Linux kernel, and device drivers, across ARM64 and x86 server and storage systems. Contributed to the design and evolution of a custom BMC firmware stack supporting large‑scale fleet management in hyperscale data‑center environments. 
          \item Composable Storage and High‑Speed Networking Engineering. Architected and developed a composable storage platform leveraging PCIe and Ethernet‑based fabric backbones, enabling dynamic resource pooling and high‑bandwidth, low‑latency data paths across distributed systems. Implemented and validated firmware for 100G networking subsystems, including PHY bring‑up, link training, diagnostics, and performance tuning across multi‑lane high‑speed SerDes interfaces. Ported Broadcom’s 32-bit Network Operating System to an ARM64 hardware platform (Applied Micro Xgene), integrating: Bootloader and BSP modifications Hardware abstraction layer (HAL) enablement, Platform drivers for networking, storage, and management subsystems, System bring‑up, debugging, and performance validation. Collaborated closely with hardware, ASIC, and platform teams to ensure seamless interoperability across PCIe, Ethernet, and ARM64 components..
          \item Developed UEFI software applications and drivers for data‑center storage platforms, enabling secure firmware update management and comprehensive pre‑boot system health checks, including MCA, RAS, and DIMM ECC error reporting. Collaborated with third‑party silicon and device vendors to design and integrate UEFI device drivers, ensuring reliable firmware update workflows and advanced diagnostic capabilities for external SoCs. Contributed to the successful launch of multiple products by delivering robust firmware components and driving effective cross‑functional collaboration across hardware, software, and platform teams.
          \item Developed a proof‑of‑concept validation platform for the next‑generation InfiniBand fabric backend powering EMC VMAX storage systems. Designed and implemented a bare‑metal C driver for the MLX4 SoC and built a modular C/Python test framework capable of orchestrating complex, high‑stress performance workloads. Created custom MLX MAD packet generators to characterize bandwidth, latency, and reliability limits across the fabric. The platform became a critical tool for early architectural decision‑making—exposing integration defects, uncovering key performance behaviors, and directly shaping the design of the new fabric subsystem. This work substantially accelerated the engineering team’s ability to test, iterate, and converge on a robust system architecture.
          \item Developed a family of custom ASICs integrating PCIe–sRIO interfaces for EMC Symmetrix enterprise data‑storage systems. Led the architecture and design of custom Tensilica‑based embedded processors, including instruction‑set extensions and firmware ecosystem development. Served as Technical Leader for firmware architecture, guiding system‑level design decisions, development processes, and cross‑functional integration. Designed and implemented embedded C firmware optimized for custom Tensilica processor cores. Created hardware designs using Verilog and Cadence toolchains, ensuring tight coupling between hardware and firmware components. Drove hardware–firmware co‑design, performing tradeoff analysis and optimization to meet performance, power, and area targets.
      \end{itemize}}


      \resumeExpSubheading
      {EMC Corporation}{Cork, Ireland}
      {\makecell[l]{ASIC Component Engineer \\ Test Engineer}}{\makecell[r]{June 1997 - Dec. 1999 \\ June 1993 - Aug. 1997}}
      {\skill{HW Debugging} \skill{Data Storage} \skill{Testing} \skill{ASIC/FPGA Technology}}
      \resumeDesc{\begin{itemize}
          \item Built and led an internal environmental test organization responsible for stress‑testing Symmetrix controller boards to validate reliability and performance under extreme thermal, electrical, and mechanical conditions.
            \item Designed and commissioned specialized environmental and diagnostic test equipment, created end‑to‑end test strategies, and delivered training programs for engineering and manufacturing teams to ensure consistent execution.
            \item Performed advanced board‑ and component‑level debugging on Symmetrix enterprise storage systems, including deep‑dive ASIC and FPGA failure analysis, to identify root causes, implement corrective actions, and improve product robustness..
      \end{itemize}}


      \resumeExpSubheading
      {Amdahl Corporation}{Dublin, Ireland}
      {Test Engineer}{June 1992 - June 1993 }      
      {\skill{HW Debugging} \skill{Mainframe computer architecture}}
      \resumeDesc{\begin{itemize}
          \item Assigned to Test Engineer role for the Amdahl 5995M mainframe. \item Responsible for board level test and debug of various elements within the mainframe.
      \end{itemize}}
%
      
      \resumeExpSubheading
      {Cork Institute of Technology}{Cork, Ireland}
      {Research Assistant}{Jan. 1992 - June 1992 }      
      {\skill{HW,SW Design} \skill{System Requirements} \skill{8051 assembly}}
      \resumeDesc{\begin{itemize}
          \item Designed and implemented a “Talking Multi-meter”.
          \item Responsible for the design of an 8051 Embedded Microcontroller and a TI speech synthesis chip for taking electrical measurements as input and used synthetic speech to output the result of the measurement.
          \item The project included both hardware design and software development using the 8051 assembly language.
      \end{itemize}}
      
      \resumeSubHeadingListEnd
%
%-----------EDUCATION-----------------
\section{Education}
  \resumeSubHeadingListStart
    \resumeSubheading
       {University of Limerick}{Limerick, Ireland}
      {Master's degree in Computer Engineering (Part-time)}{Sept. 1997 - June. 2001}
  \resumeSubheading
       {British Computer Society / Dublin Institute of Technology (DIT)}{Dublin, Ireland}{Degree Computer Science (Part-time)}{Sept. 1992 - June. 1996}
    \resumeSubheading
      {Cork Institute of Technology CIT)}{Cork, Ireland}
      {Electronic Engineering}{Sept. 1988 - June. 1991}
	 %\begin{comment}
	 \resumeItemListStart
        \resumeItem{}
          {}
      \resumeItemListEnd
      %\end{comment}
  \resumeSubHeadingListEnd

%
%--------PATENTS------------
\section{Patents}
 \resumeSubHeadingListStart
   \item{
     \textbf{7987229}{: \href{https://patents.google.com/patent/US7987229B1/en?oq=7987229}{Data storage system having plural data pipes}}
   }\vspace{-7pt}
   \item{
     \textbf{8156220}{: \href{https://patents.google.com/patent/US8156220B1/en?oq=8156220} {Data Storage System}}
   }\vspace{-7pt}
   \item{
     \textbf{8090789}{: \href{https://patents.google.com/patent/US8090789B1/en?oq=8090789} {Method of operating a data storage system having plural data pipes}}
   }\vspace{-7pt}
   \item{
     \textbf{8645623}{: \href{https://patents.google.com/patent/US8645623B1/en?oq=8645623} {Method for performing a raid operation in a data storage system}}
   }
   \vspace{-7pt}
   \item{
     \textbf{7631128}{: \href{https://patents.google.com/patent/US7631128B1/en?oq=7631128} {Protocol controller for a data storage system}}
   }
 \resumeSubHeadingListEnd


%-------------------------------------------
\end{document}

